\documentclass[a4paper,11pt]{article}
\usepackage{jheppub} % for details on the use of the package, please see the JINST-author-manual
\usepackage{lineno}
% \linenumbers
\newcommand{\rmd}{\mathrm{d}}
\usepackage{eufrak}
\usepackage{amsmath}
\usepackage{mathtools}
\newcommand{\mynote}[1]{{\textcolor{red}{[#1]}}}
\newcommand{\Anote}[1]{{\textcolor{magenta}{[#1]}}}

% \arxivnumber{2304.XXXXX} % if you have one

\title{\boldmath Critical JT Gravity}

% Collaborations

%% [A] If main author
%% \collaboration{\includegraphics[height=17mm]{collabroation-logo}\\[6pt]
%%  XXX collaboration}

%% or
%% [B] If "on behalf of"
%% \collaboration[c]{on behalf of XXX collaboration}


% Authors
% The "\note" macro will give a warning: "Ignoring empty anchor...", you can safely ignore it.

%% [A] simple case: 2 authors, same institution
%% \author[1]{A. Uthor\note{Corresponding author.}}
%% \author{and A. Nother Author}
%% \affiliation{Institution,\\Address, Country}

%% or, e.g.
%% [B] more complex case: 4 authors, 3 institutions, 2 footnotes
%% \author[a,b]{F. Irst,\note{Now at another university}}
%% \author[c]{S. Econd,}
%% \author[a,2]{T. Hird\note{Also at Some University.}}
%% \author[c,2]{and Fourth}
%% \affiliation[a]{Institution_1,\\Address, Country}
%% \affiliation[b]{Institution_2,\\Address, Country}
%% \affiliation[c]{Institution_3,\\Address, Country}

\author{Alicia Castro}
\affiliation{IMAPP, Radboud University,\\
6525 AJ Nijmegen, Netherlands}
% \affiliation{Another University,\\
% different-address, Country}

% E-mail addresses: only for the corresponding author
\emailAdd{alicia.castro@science.ru.nl}
\abstract{In this paper, we investigate a critical behavior of JT gravity, a model of two-dimensional quantum gravity on constant negatively curved spacetimes. Our approach involves using techniques from random maps to investigate the generating function of Weil--Petersson volumes, which count random hyperbolic surfaces with defects. The defects are weighted geodesic boundaries, and criticality is reached by tuning the weights to the regime where macroscopic holes emerge in the hyperbolic surface, namely \textit{non-generic criticality}. We analyze the impact of this critical regime on some universal features, such as its density of states. We present a family of models that interpolates between systems with $\rho_0(E)\sim\sqrt{E-E_0}$ and $\rho_0(E)\sim (E-E_0)^{3/2}$, which are commonly found in models of JT gravity coupled to dynamical end-of-the-world and FZZT branes, and give a precise definition of what this phase transition means from the random geometry point of view.}

\begin{document}
\maketitle
\flushbottom
% yes, it's the generating fuction of trees with vertex weight $\mu$\\
% combinatorial explanation of x(r)=r\\
% doesn't generate anything, the solution is R=0

\section{Introduction}

Quantizing gravity remains one of the most challenging and unresolved problems in physics. The traditional methods for quantizing physical systems fail to apply when it comes to gravity. For instance, when using the path integral formalism, one needs to integrate over the space of all possible 4-dimensional spacetimes governed by Einstein-Hilbert's action. However, this space is not yet completely understood mathematically, nor are the physical constraints, such as topology change and Lorentzian versus Euclidean signatures \cite{Ambjorn:1998fd}, well-defined.  As a result, studying lower-dimensional toy models is often a useful tool for analyzing these issues.

Two-dimensional quantum gravity is a useful toy model from both the physical and mathematical perspectives. Physically, 2D quantum gravity allows us to study key features of gravity without having to deal with propagating degrees of freedom i.e.\ gravitons  Mathematically, the existence of conformal symmetry and discretization techniques in two dimensions make some aspects of 2D quantum gravity computable. In particular, the 2D Einstein-Hilbert action is topological, so the euclidean path integral reduces to counting geometries with a weight given by $e^{-S_0\chi}$, where $\chi$ is the Euler characteristic and $S_0\sim\frac{1}{G}$ (see \cite{Budd:2022zry} for an updated review). Yet, the exact quantization of 2D gravity remains a compelling task.

Recent developments in JT gravity \cite{Jackiw:1984je,Teitelboim:1983ux}, a model of euclidean 2-dimensional quantum gravity on constant negatively curved spacetimes (the hyperbolic plane, which is the euclidean version of AdS$_2$), have proved potential for shedding light on this topic due to its simplicity and rich physical content \cite{Mertens:2022irh}. One of the fascinating features of JT gravity is its connection to random geometries \cite{Witten:2020bvl}.  In particular, it has been demonstrated that the partition function of JT gravity can be split into two parts: one is proportional to volumes of moduli spaces of hyperbolic surfaces, while the other only involves information about the asymptotic behavior of AdS$_2$. The latter component can be interpreted in terms of a holographic correspondence, as it is the partition function of a quantum mechanical system given by the Schwarzian derivative effective action near the asymptotic boundary \cite{Maldacena:2016upp, Kitaev}. Furthermore, the Weil-Petersson volumes carry the geometrical information of the negatively-curved surface. In addition, studying the higher genus contributions to this partition function revealed its relation with yet another random geometry field, matrix models \cite{Saad:2019lba,Stanford:2019vob}. This allowed for significant progress in the field, since the formal machinery of matrix models \cite{Eynard:2016yaa} can be applied to JT gravity \cite{Johnson:2019eik}. 


In this paper, we apply novel techniques from random geometry \cite{budd2023} to JT gravity. We begin by examining the generating function of Weil-Petersson volumes with non-trivial weights on its geodesic boundaries. This allows us to investigate the behavior of the JT gravity partition function, in particular, the regime where there is a non-analytical behavior of the disk function and density of states, which are indicative of non-generic criticality. Such critical regimes are reminiscent from the planar map case where there is a transition from a phase of  surfaces with microscopic holes to a phase with macroscopic ones (see Figure \ref{fig:JT_picture_holes}). This translates, in the hyperbolic geometry case, to the proliferation of large geodesic boundaries. This phase transition is driven by a fine-tuning of the geodesic weights and has a significant impact on the density of states of the theory in the boundary. Our approach allows us to identify a family of models that interpolates between systems with $\rho_0(E)\sim\sqrt{E}$ and $\rho_0(E)\sim E^{3/2}$, which are commonly found in JT gravity coupled to dynamical end-of-the-world (EOW)\cite{Gao:2021uro} and FZZT branes \cite{Okuyama:2021Page}.\\
\begin{figure}[htbp]
\centering
\includegraphics[width=.5\textwidth]{JT_withholes.pdf}
\caption{Example of a configuration of critical JT gravity. The blue line represents the nearly AdS boundary and in red are the geodesics boundaries. In the non-generic critical regime, these boundaries become macroscopic.\label{fig:JT_picture_holes}}
\end{figure}\\
This paper is organised as follows. In section \ref{section:Random_geometry} we introduce random maps and random hyperbolic surfaces as examples of random geometries, together with the notion of criticality. In section \ref{section_JT_random}, we introduce a dictionary between the vocabularies of random geometry and JT gravity using a common tool, the string equation. To clarify, we reproduce results on FZZT branes and conical defects in these terms. In Section \ref{Section_Criticality}, we define and analyse the different critical regimes of the string equation and we focus on its implication for the partition and disks functions and density of states in the non-generic critical phase. In section \ref{Section_Dilaton_gravity}, we construct the corresponding dilaton potential and analyse the thermodynamical properties of these spacetimes and their relation to previous studies using branes and matrix models. In Section \ref{section:SFF}, we compute the spectral form factor of this family of models and comment on its potential to reach a late-time plateau. Finally, in section \ref{Section_Discussion}, we give an outlook and discuss potential implications for JT gravity. 

\section{Random Geometry}\label{section:Random_geometry}
In order to make this paper as self-contained as possible, let us recall some terms that are useful to describe random geometries, for more details see \cite{Eynard:2016yaa} and \cite{Budd:2022zry} for an updated review. A \textbf{\textit{graph}} is a finite set of vertices, edges and an incidence map  $\left\{ v,e,\mathbb{I}\right\} $, $\mathbb{I}:e\rightarrow (v_1,v_2)$ that associates two vertices with an edge. A \textbf{\textit{genus-g map}} is an equivalence class of \textit{graphs} embedded in a surface of genus-g, such that its edges do not cross, modulo orientation-preserving homeomorphisms of the surface.  In particular, throughout this paper, when we talk about genus-$g$ maps, we will only consider maps with one marked edge (called the root), a defined orientation (an arrow that goes from one vertex to the other on the root edge), pointed (one distinguished vertex of the root) and an even number of edges per face. In the case $g=0$, these are called rooted pointed planar bipartite maps. A face of a genus-g map is a topological disk bounded by edges. The \textit{length of a face} is the sum of the lengths of its edges, where each edge has length $1$.

%A \textbf{\textit{Boltzmann genus-g map}} is a \textit{genus-g map} with admissible non-negative values on the faces, which we call \textit{weights}.\mynote{This does not really make sense: the Boltzmann genus-g map is a random genus-$g$ map, whose distribution is determined by certain weights. It is not a map with values on the faces (for starters the values can only depend on the length/degree)} %%These weights can be understood as a decoration of the metric map which can be associated with a coupling to matter.

\subsection{Random maps}
Let $\mathcal{M}_{g,n}(L_1,\dots,L_n)$ be the set of genus-$g$ maps with $n$ faces of half-lengths\footnote{Recall, all faces have an even number of edges, so their length is $2L_i$.} $\{L_i\in \mathbb{N}\}_{1\leq i\leq n}$ and weights $\mathbf{q}=(q(L_i)\geq 0)$ per face of half-length $L_i$ and with $q(0)$ is the weight per vertex which we will usually set to $1$. These weights can be understood as a decoration of the map which can be associated with a coupling to matter. The cardinality of this set is the total number of Boltzmann genus-$g$ maps up to orientation-preserving automorphisms that preserve the face labels
\begin{equation}
    ||\mathcal{M}_{g,n}(L_1,\dots,L_n)||=\sum_{\mathfrak{m}\in\mathcal{M}_{g,n}({L_i})}\frac{1}{\mathrm{Aut}(\mathfrak{m})}.
\end{equation}
For fixed values of $g$ and $n$, these are nothing more than positive real numbers and they can be encoded in a generating function
\begin{equation}
  F_g[\mathbf{q}]= \sum_{k=1}^\infty\frac{1}{k!} \sum_{L_1=0}^\infty q(L_1) \cdots \sum_{L_k=0}^\infty q(L_k) q(0)^{2-2g-k+\sum_i L_i} ||\mathcal{M}_{g,k}(L_1,\dots,L_k)||.
\end{equation}
where the powers of $q(0)$ account for the number of vertices with weight $q(0)$ in a map of genus-$g$ with $k$ faces of half-length $\sum_i L_i$ according to Euler's characteristic. Therefore, we can write  
\begin{equation}
     ||\mathcal{M}_{g,n}(L_1,\dots,L_n)||=q(0)^{-2+2g+n-\sum_i L_i}\left.\frac{\partial^n F_g[\mathbf{q}]}{\partial q(L_1)\dots \partial q(L_n)}\right|_{\mathbf{q}=0}.
\end{equation}
$F_g[\mathbf{q}]$ is called the \textit{partition function of genus-g maps}. It can be observed that the weights $q(L_k)$ control the likelihood of having faces of length $L_k$ in the map, e.g. if $\mathbf{q}=(q(0)=1,q(1)=0,q(2)\neq 0,q(3)=0,0,0,\dots)$, then only faces of length $2\cdot2$ are allowed and the genus-$g$ map is a quadrangulation of a genus-$g$ surface. Coming back to the generating function $F_g[\mathbf{q}]$, in order to analyse it, it is convenient to write it in terms of another power series, the \textit{string equation} 
\begin{equation}
   \tilde{x}(u)=u-\sum_{k=0}^\infty q(k)\;\frac{1}{2}\binom{2k}{k} u^k,\label{eq:string_eq_maps}
\end{equation}
where $q(k)\coloneqq q(L=k)$ with $k\in \mathbb{N}$. In particular, the genus-$0$ generating function is proven to be \cite{Tutte1968OnTE,Eynard:2016yaa}\footnote{For $g>0$, check \cite{Eynard:2016yaa}.}
\begin{align}
    F_0 &= \frac{1}{2}\int_{0}^{E_\mathbf{q}} \!\!\!\rmd u\,\frac{1}{u} \left(\tilde{x}(u)^2 -(u-q(0))^2 \mathbb{I}_{u<q(0)} \right),\label{eq:genfun0beta}
\end{align}
where $E_\mathbf{q}>0$ is the smallest solution to the string equation $\tilde{x}(E_\mathbf{q})=0$. Thus, $E_\mathbf{q}$ is an implicit function of the weight $\mathbf{q}$. The weight is said to be \textit{admissible} if such a solution exists. Furthermore, another useful quantities we are interested are the \textit{moment} $M_0$, which is defined 
\begin{equation}
M_0=\left(\frac{\partial E_\mathbf{q}}{\partial q(0)}\right)^{-1}=\frac{\partial \tilde{x}}{\partial u}(E_\mathbf{q}),
\end{equation}
% or, equivalnetly
% \begin{equation}
%    x(u)=u-\sum_{m=0}^\infty\frac{t_m}{m!}u^m,
% \end{equation}
and the disk function, which corresponds to the generating function of genus-$0$ maps with one marked face and an arbitrary number of unmarked faces, operationally, this is 
\begin{equation}
    W_\mathbf{q}(L) = \frac{\partial 	F_0[\mathbf{q}]}{\partial q(L)}.
\end{equation}
It can also be noted that
$\frac{\partial^2	F_0[\mathbf{q}]}{\partial q(0)\partial q(1)}=E_\mathbf{q}
$
which gives a geometrical meaning to $E_\mathbf{q}$, it is the generating function of rooted, pointed planar maps.\\
By this point, all our expressions are formal power series, but in order to talk about random geometry and random maps, we need a well-defined probability measure. This can be also understood as a well-defined canonical way of sampling a map uniformly. In particular, we would like to define a probability density measure on the space of planar maps with $k$ faces with weight $\mathbf{q}$. This can be done using the generating function in the following way   
\begin{equation}
    p\left(L_1, ... , L_k\right) = \frac{1}{k!} \frac{\prod_{i=1}^{k} q(L_i)}{F_0[\mathbf{q}]}.
\label{eq:prob_density_boltzmann}
\end{equation}
Note that this is only well defined when $\mathbf{q}$ is such that $F_0[\mathbf{q}]<\infty$. This is satisfied iff the weight $\mathbf{q}$ is admissible, i.e.\, if the string equation has a solution $E_\mathbf{q}>0$. In this way, given an admissible weight, a \textbf{\textit{Boltzmann planar map}} is defined as a random planar map with $k$-faces of random lengths sampled with probability \eqref{eq:prob_density_boltzmann}.
\subsection{Critical maps and the $O(n)$ loop model}
Now that we introduced Boltzmann planar maps, we can study more of their characteristics. In particular, the main point we want to make is that the string equation is the central building block to construct all we have done so far. Consider the solution to the string equation $\tilde{x}(E_\mathbf{q})/E_\mathbf{q}=0$ and its first derivative $\tilde{x}'(E_\mathbf{q})$,
\begin{equation}
    1-E_\mathbf{q}^{-1}=\sum_{k=1}^\infty q(k)\;\frac{1}{2}\binom{2k}{k} E_\mathbf{q}^{k-1}, \hspace{1cm} \tilde{x}'(E_\mathbf{q})=1-\sum_{k=1}^\infty q(k)\;\frac{k}{2}\binom{2k}{k} E_\mathbf{q}^{k-1}.
\end{equation}
Given that $q(k),E_\mathbf{q}>0$, we can conclude that $0\leq \tilde{x}'(E_\mathbf{q})<1$.
In the mathematics literature, if an admissible weight is such that $0<\tilde{x}'(E_\mathbf{q})$, it is called \textit{subcritical}, and if $\tilde{x}'(E_\mathbf{q})=0$, the weight is called \textbf{\textit{critical}}. Furthermore, if a weight is critical, it can be \textit{generic critical} if the string equation behaves as $\tilde{x}(u) \sim (E_\mathbf{q}-u)^2$ for $u\lesssim E_\mathbf{q}$ and \textbf{\textit{non-generic critical}} if the string equation behaves as $\tilde{x}(u)\sim (E_\mathbf{q}-u)^{\alpha-1/2}$ with $3/2<\alpha<5/2$ for $u\lesssim E_\mathbf{q}$ \cite{10.2307/25450028}. In this last case, $E_\mathbf{q}>0$ is the radius of convergence of the formal power series $\tilde{x}(u)$.
\begin{table}[h]
    \centering
    \begin{tabular}{|c|c|}
    \hline
        Weight $\mathbf{q}$ & String equation \\
    \hline     
         Subcritical & $\tilde{x}'(E_\mathbf{q})> 0$ 
         \\
         Generic critical & $\tilde{x}'(E_\mathbf{q})=0, \hspace{2mm} \tilde{x}(u) \sim (E_\mathbf{q}-u)^2$
         \\
         Non-generic critical & $\tilde{x}'(E_\mathbf{q})=0, \hspace{2mm} \tilde{x}(u)\sim (E_\mathbf{q}-u)^{\alpha-1/2}$\\
         \hline
    \end{tabular}
    \caption{Types of admissible weights.}
    \label{tab:critical}
\end{table}\\
% We are particularly interested in the scaling limit of maps with non-generic critical weights, also called \textbf{\textit{stable maps}}.
In the non-generic critical case, it is proven that the weight must behave as $\mathbf{q}(k)\sim h_1^{k}k^{-\alpha}$ for some $h_1>0$ and $k\rightarrow\infty$ \cite{legall,buddcurien}. The geometrical meaning of this slow decay of critical weights is the proliferation of large faces $k\gg 1$ (See right part of Figure \ref{fig:holes_maps}). 
To further understand this, we will follow the reasoning used to draw a relation between these maps and the $O(n)$ loop model in \cite{Borot:2011kg,2018arXiv180902012B}. In that case, one considers a planar map dressed with a loop configuration which is
a collection of disjoint loops drawn on the planar map each of arbitrary length $2k$, where the parameter $n$ controls the weight per loop. The external gasket $E(n,\Gamma)$ is then obtained by removing all edges and vertices situated within the interior of the loops in $\Gamma$, as well as those intersected by them, from the planar map, these faces of the exterior gasket are called \textit{holes} (see Figure \ref{fig:holes_maps}). The random gasket configurations $O(n)$ loop configurations with $n=2 \sin\left((\pi(\alpha-1/2)\right)$ are related to non-generic critical Boltzmann planar maps for $\alpha\in (3/2,5/2)$ \cite{Borot_2012}. The construction consists on substituting the faces touched and contained on each loop by holes with weights assigned consistently with the disk functions relation $W^{(n)}_\mathbf{q}(k)=W_\mathbf{\hat{q}}(k)$, where $W^{(n)}_\mathbf{q}(k)$ is the disk function of the $ O(n)$ loop model with weights $\mathbf{q}$, and $\hat{\mathbf{q}}(k)=\mathbf{q}(k)+n\ h_1^{k} W_\mathbf{\hat{q}}(k)$. We will make this clear by studying this explicit construction in the hyperbolic case.
\begin{figure}[htbp]
\centering
\includegraphics[width=.8\textwidth]{boltzmann_holes_2.pdf}
\caption{\textbf{Left}: Illustrative example of a $O(n)$ model configuration with two non-intersecting loops in red. The interiors of the loops are colored grey. \textbf{Right}: The corresponding external gasket where the marked faces, or holes, are colored in gray.\label{fig:holes_maps}}
\end{figure}\\
We will come back to this topic in the next subsection and elaborate more on the hyperbolic surfaces case, but interested readers can refer to \cite{Borot:2011kg} for proofs concerning the $O(n)$ loop model in the context of maps.

% One can then prove the asymptotic behavior of various quantities, including the disk function 
% \begin{equation}
%     W(L) \sim t \; e^{L\sqrt{E_\mathbf{q}}}L^{-\alpha-1}\label{disk_critical}.
% \end{equation}

%%%%%%%%%%%%%%%%%%%%%%%%%%%%%%%%%%%%%%

\subsection{Random hyperbolic surfaces}
Along the lines of random planar maps, we want to illustrate how to consistently define a random hyperbolic surface by constructing a probability measure in the space of hyperbolic surfaces of genus $g$ and $k$ geodesic boundaries, $\mathcal{M}^{WP}_{g,k}(b_1,\dots,b_k)$. The starting point is to write the `cardinality' of this space, this is given by the volume of the moduli space of $\mathcal{M}^{WP}_{g,k}(b_1,\dots,b_k)$, called Weil-Petersson volumes $V^{WP}_{g,k}(b_1,\dots,b_k)$ \cite{Mirzakhani}. The second step is to write a generating function of Weil-Petersson volumes of genus $g$ \cite{1996CMaPh.181..763K}
\begin{align}
	F^{WP}_g[\mu] &= \sum_{k=0}^\infty\frac{2^{2-2g-k}}{k!} \int_{0}^\infty \rmd \alpha_1 \mu(\alpha_1) \cdots \int_{0}^\infty \rmd \alpha_k \mu(\alpha_k) V^{WP}_{g,k}(\alpha_{1}\ldots,\alpha_{k}),
\end{align}
Note that similarly to the Boltzmann map case, $\mu(\ell_k)$ can be interpreted as the likelihood of having geodesic boundaries of length $\ell_k$, i.e.\ it is a weight\footnote{Another way to see this is to consider a weight $\mu(\ell)=x\tilde{\mu}(\ell)$ to make a clear distinction between the formal variable of the generating function, $x$, and the statistical weight $\tilde{\mu}$ in the sense of a non-negative valued label on the geodesic boundaries of the hyperbolic surface.}. In Section \ref{section_JT_random} we will explain the role that weights play in JT gravity with defects.\\
The generating function of Weil-Petersson volumes can be expressed in terms of the string equation 
\begin{equation}
   x(u)=u-\sum_{k=0}^\infty\frac{t_k+\gamma_k}{k!}u^k,\label{eq:string_eq_times}
\end{equation}
where 
\begin{equation}
\gamma_k=\frac{(-1)^k\pi^{2k-2}}{(k-1)!}\mathrm{1}_{k\geq 2}\label{gamma_k}
\end{equation}
are the KdV times \cite{1996CMaPh.181..763K, Mirzakhani} and are called the times of pure JT gravity in the physics literature \cite{Johnson:2019eik,Okuyama:2019xbv}, and
\begin{equation}
    t_k=\frac{4^{-k}}{k!}\int_0^\infty \rmd \ell \mu(\ell) \ell^{2k}\label{eq:def_general_JT_times}
\end{equation}
are the times associated with geodesic boundaries with labels $\mu(\ell)$. Note that for Boltzmann maps, we have $\gamma_k=0$.
 Equivalently, the string equation \eqref{eq:string_eq_times} can be written in the following form 
\begin{equation}
    x(u)=\frac{\sqrt{u}}{\pi}J_1\left(2\pi\sqrt{u}\right)-\int_0^\infty \rmd\ell \mu(\ell) I_0\left(\ell \sqrt{u}\right).\label{eq:general_string_eq_JT}
\end{equation}
The first term is well-known in the physics literature as \textit{pure JT gravity}, up to a change of sign $u\rightarrow -u$. The second term accounts for the presence of geodesic boundaries of weight $\mu(\ell)$ and it has been used in the physics literature couple JT gravity to conical defects and branes \cite{Turi:2020ale,Okuyama:2021FZZT,Gao:2021uro}.\\
In the rest of this section, we will restrict our study to the case $g=0$ for convenience, but the interested reader can check \cite{budd2022irreducible,budd2023_2} for $g>0$. The genus-$0$ generating function of Weil-Petersson volumes can be expressed in terms of the string equation in the same way as for Boltzmann planar maps, as well as the moment $M_0$ \cite{budd2022irreducible}
\begin{align}
    	F^{WP}_0[\mu] &= \frac{1}{4}\int_{0}^{E_\mu} \!\!\!\rmd u\, x(u)^2,\label{eq:genfun0}\\
    M^{WP}_0&=\frac{\partial x}{\partial u} (E_\mu)\label{eq:m0_wp}
\end{align}
where $E_\mu>0$ is the smallest solution to $x(E_\mu)=0$, and, in the physics literature, it is called the threshold energy $E_\mu=-E_0$.  This allows us to define a probability density measure on the space of hyperbolic genus-$0$ surfaces  
\begin{equation}
    p^{WP}\left(\ell_1, ... , \ell_k\right) = \frac{1}{k!} \frac{\prod_{i=1}^{k}  \mu(\ell_i)}{F^{WP}_0[\mu]}.
\label{eq:prob_density_WP}
\end{equation}
In this way, given an admissible weight $\mu$ such that $F^{WP}_0[\mu]<\infty$, we can call \textbf{\textit{random genus-$0$ hyperbolic surface with $k$-boundaries of random lengths}} to a surface sampled with probability \eqref{eq:prob_density_WP}.\\
In particular, once again, it is useful to introduce the disk function, this is the generating function of hyperbolic surfaces with one marked geodesic boundary of length $b$
\begin{align}
W_\mu(b) &=\frac{\delta F^{WP}
_0}{\delta\mu(b)}=\sum_{k=2}^\infty\frac{2^{2-(k+1)}}{k!}\int_0^\infty \rmd \ell_1
\mu(\ell_1)...\int_0^\infty \rmd \ell_k
\mu(\ell_k)V^{WP}_{0,k+1}(b,\ell_1,...,\ell_k)\\
   &= -\frac{1}{2}\int_{0}^{E_\mu} \!\!\!\rmd u\, x(u)\;I_0(\ell\sqrt{u}).\label{eq:JT_disk_function}
\end{align}
%%%%%%%%%%%%%%%%%%%%%%%%%%%%
\subsection{Critical hyperbolic surfaces and the hyperbolic $O(n)$ loop model}
A natural question that arises is if a critical behavior, analogous to that observed in planar maps, happens for hyperbolic surfaces too. Specifically, whether there exist weights $\mu$ such that the string equation \eqref{eq:general_string_eq_JT} behaves in a similar manner to the ones presented in Table \ref{tab:critical}. The key difference between the Boltzmann maps and hyperbolic surfaces string equations are the background times $\gamma_k$. This implies that, in order to reach a non-generic critical regime in the hyperbolic case, the weight $\mu$ needs to be fine-tuned to dominate the Bessel function in the first term of \eqref{eq:general_string_eq_JT}, instead of the simpler linear term in the Boltzmann map case \eqref{eq:string_eq_maps}. Therefore, it appears to be a significantly larger challenge to find non-generic critical weights for hyperbolic surfaces. Within this subsection, we will establish their existence by drawing upon an analogy with the gasket decomposition of the hyperbolic $O(n)$ loop model, while interested readers can refer to \cite{budd2023} for comprehensive formal proofs.\\
Consider the string equation \eqref{eq:general_string_eq_JT}, which, for convenience, we write again 
\begin{equation*}
 x(u)=\frac{\sqrt{u}}{\pi}J_1\left(2\pi\sqrt{u}\right)-\int_0^\infty \rmd\ell \mu(\ell) I_0\left(\ell \sqrt{u}\right),    
\end{equation*}
and a weight $\mu(\ell) = \kappa \,e^{-\lambda\ell} \hat{\mu}(\ell)$ with $\kappa,\lambda>0$ and $\hat{\mu}(\ell)$ a bounded, non-negative function on $[0,\infty)$ with asymptotics $\hat{\mu}(\ell) \sim C \,\ell^{-\alpha}$ as $\ell\to\infty$, e.g.\ $\hat{\mu}(\ell) = (\ell+1)^{-\alpha}$ or $\hat{\mu}(\ell) = \ell^{-\alpha} \Theta(\ell-1)$. We will demonstrate how we can choose constants $\kappa,\lambda>0$ to reach non-generic criticality, i.e.\ $x(E_\mu)=0$, $x'(E_\mu)=0$ and $x(u) \sim c\, (E_\mu-u)^{\alpha-1/2}$ as $u\to E_\mu$, with $\alpha \in (3/2,5/2)$.
To see this, we note that if we have $\mu$ such that $x'(E_\mu)=0$, then  $x''(E_\mu)=\infty$. This last expression requires that $\lambda = \sqrt{E_\mu}$.
Then, $x(E_\mu)=0$ allows us to fix $\kappa$ to
\begin{equation}
    \kappa = \frac{\frac{\sqrt{E_\mu}}{\pi} J_1(2\pi\sqrt{E_\mu})}{\int_0^\infty \rmd \ell \hat{\mu}(\ell) e^{-\ell\sqrt{E_\mu}}I_0(\ell \sqrt{E_\mu})}.\label{eq:kappa}
\end{equation}
If we take $E_\mu$ to be the first zero of
\begin{equation}
    x'(E_\mu) = J_0(2\pi\sqrt{E_\mu}) - \frac{\sqrt{E_\mu}}{\pi} J_1(2\pi\sqrt{E_\mu}) \frac{\int_0^\infty \rmd \ell \hat{\mu}(\ell) e^{-\ell\sqrt{E_\mu}}\frac{\ell I_1(\ell \sqrt{E_\mu})}{\sqrt{E_\mu}}}{\int_0^\infty \rmd \ell \hat{\mu}(\ell) e^{-\ell\sqrt{E_\mu}}I_0(\ell \sqrt{E_\mu})},
\end{equation}
which exists since $x'(0)=1$ and it is negative for $\sqrt{2}E_\mu = c_0^2/(8\pi^2) = 0.0732\ldots$, where $c_0$ is the first zero of the Bessel function $J_0$.
Accordingly,
%\begin{align}
%    x''(u) &= \left(\frac{\sqrt{u}}{\pi}J_1\left(2\pi\sqrt{u}\right)\right)'' - \int_0^\infty \rmd \ell \mu(\ell) \frac{\ell^2}{u}I_2(\ell\sqrt{u})\\
 %   &= (\cdots)''-\int_0^\infty \rmd \ell \ \kappa \ C \,e^{-\sqrt{E_\mu}\ell} \ell^{-\alpha} \frac{\ell^2}{u}\frac{e^{\sqrt{E_\mu}\ell}}{\sqrt{2\pi\sqrt{E_\mu}\ell}}+\mathcal{O}(\ell^{-\alpha+5/2})\\
  %  &= (\cdots)''-\frac{\kappa \ C}{u\sqrt{2\pi\sqrt{E_\mu}}} \frac{1}{\frac{5}{2}-\alpha}\ \ell^{-\alpha+5/2} + \mathcal{O}(\ell^{-\alpha+5/2})\\
   % &= -\frac{\kappa \ C (\frac{3}{2}-\alpha)}{\sqrt{2\pi\sqrt{E_\mu}}} \Gamma\left(\frac{3}{2}-\alpha\right)\ (E_\mu-u)^{\alpha - 5/2} + \mathcal{O}((E_\mu-u)^{\alpha - 5/2}),
%\end{align}

\begin{align}
    x''(u) &= \left(\frac{\sqrt{u}}{\pi}J_1\left(2\pi\sqrt{u}\right)\right)'' - \int_0^\infty \rmd \ell \mu(\ell) \frac{\ell^2}{u}I_2(\ell\sqrt{u})\\
    &= (\cdots)''-\int_0^\infty \rmd \ell \ \kappa \ C \,e^{-\sqrt{E_\mu}\ell} \ell^{-\alpha} \frac{\ell^2}{u}\frac{e^{\sqrt{u}\ell}}{\sqrt{2\pi\sqrt{u}\ell}}+\mathcal{O}(\ell^{-\alpha+5/2})\\
    &= (\cdots)''-\frac{\kappa \ C}{u\sqrt{2\pi\sqrt{u}}} \Gamma\left(\frac{5}{2}-\alpha\right)\ (\sqrt{E_\mu}-\sqrt{u})^{\alpha - 5/2} + \mathcal{O}((\sqrt{E_\mu}-\sqrt{u})^{\alpha - 5/2})\\
    &= -\frac{\kappa \ C}{E_\mu \sqrt{2\pi\sqrt{E_\mu}}} \Gamma\left(\frac{5}{2}-\alpha\right)\ (2\sqrt{E_\mu})^{5/2-\alpha}(E_\mu-u)^{\alpha - 5/2} + \mathcal{O}((E_\mu-u)^{\alpha - 5/2}),
\end{align}
where we use Karamata's Tauberian Theorem (see e.g. \cite{feller-vol-2}, Section XIII.5, Theorem 2) in the last line. Integrating twice and using that $x(E_\mu)=x'(E_\mu)=0$ gives

%\begin{equation}
 %   x(u)= \frac{\kappa \ C}{\sqrt{2\pi\sqrt{E_\mu}}}\frac{\Gamma\left(\frac{3}{2}-\alpha\right)}{(\alpha-\frac{1}{2})}  (E_\mu-u)^{\alpha-1/2}+\mathcal{O}((E_\mu-u)^{\alpha - 1/2}).
%\end{equation}
\begin{align}
    x(u)&= -\frac{\kappa \ C}{\sqrt{2\pi}}\Gamma\left(\frac{1}{2}-\alpha\right) 2^{5/2-\alpha} E_\mu^{-\alpha/2} (E_\mu-u)^{\alpha-1/2}+\mathcal{O}((E_\mu-u)^{\alpha - 1/2})\\
    &=-\frac{\kappa \ C}{\sqrt{2\pi\sqrt{E_\mu}}}\Gamma\left(\frac{1}{2}-\alpha\right) 2^{2} (\sqrt{E_\mu}-\sqrt{u})^{\alpha-1/2}+\mathcal{O}((E_\mu-u)^{\alpha - 1/2})
\end{align}
This ends the demonstration. A similar reasoning follows for the disk function \eqref{eq:JT_disk_function}, 
\begin{align}
    W_\mu(\ell) &= -\frac{1}{2}\int_0^{E_\mu} \mathrm{d}u \,x(u) I_0(\ell\sqrt{u})\\
    &= \frac{1}{2} \frac{\kappa \ C}{\sqrt{2\pi\sqrt{E_\mu}}}\Gamma\left(\frac{1}{2}-\alpha\right) 2^{2} \int_0^{E_\mu} (\sqrt{E_\mu}-\sqrt{u})^{\alpha-1/2} \frac{e^{\ell \sqrt{u}}}{\sqrt{2\pi \ell\sqrt{u}}} \mathrm{d}u+ \cdots\\
    &=  \frac{\kappa \ C}{\sqrt{2\pi\sqrt{E_\mu}}}\Gamma\left(\frac{1}{2}-\alpha\right) 2^{2} \int_0^{\sqrt{E_\mu}} y^{\alpha-1/2} \frac{e^{-\ell y +\ell \sqrt{E_\mu}}}{\sqrt{2\pi \ell}}E_\mu^{1/4}\mathrm{d}y+ \cdots\\
 %   &= \frac{\kappa \ C}{\sqrt{2\pi\sqrt{E_\mu}}}\Gamma\left(\frac{1}{2}-\alpha\right) 2^{2} \Gamma(\alpha + 1/2) \ell^{-\alpha-1}\frac{e^{\ell \sqrt{E_\mu}}}{\sqrt{2\pi}}E_\mu^{1/4} + \mathcal{O}(\ell^{-\alpha-1})\\
    &= \frac{2}{\pi}\kappa \, C \, \Gamma\left(\frac{1}{2}-\alpha\right) \Gamma(\alpha + 1/2) \ell^{-\alpha-1}e^{\ell \sqrt{E_\mu}} + \mathcal{O}(\ell^{-\alpha-1})\\
     &= \frac{2}{\sin\left(\pi\left(\alpha-\frac{1}{2}\right)\right)}\kappa \, C \,  \ell^{-\alpha-1}e^{\ell \sqrt{E_\mu}} + \mathcal{O}(\ell^{-\alpha-1})\label{eq:disk_non-gen-crit}
\end{align}
%\begin{align}
 %   W_\mu(\ell) &= \frac{1}{2}\frac{\kappa \ C \ \Gamma\left(\frac{3}{2}-\alpha\right)  }{\sqrt{2\pi\sqrt{E_\mu}}(\alpha-\frac{1}{2})} \left(\alpha-\frac{1}{2}\right)\Gamma \left(\alpha -\frac{1}{2}\right)
  %  \, \frac{ (2 \ E_\mu)^{\alpha +\frac{1}{2}} }{\left(\ell\sqrt{E_\mu}\right)^{\alpha+\frac{1}{2}}} I_{\alpha+1/2} \left(\ell\sqrt{E_\mu}\right)+\mathcal{O}(\ell^{-\alpha-1})\\
   % &=\kappa \ C
    %\, 2^{\alpha+\frac{1}{2}} (\sqrt{E_\mu})^{\alpha-\frac{1}{2}} \frac{ \ e^{\ell \sqrt{E_\mu}}}{2\;\sin\left(\pi\left(\alpha-\frac{1}{2}\right)\right)}\ell^{-\alpha-1} + \mathcal{O}(\ell^{-\alpha-1})\label{eq:disk_non-gen-crit}
%\end{align}
where we used Euler's reflection formula for the gamma function to write the second line.
Now, we would like to show how this fits into the hyperbolic $O(n)$ loop model in a similar way as for Boltzmann planar maps \cite{Borot:2011kg}. The hyperbolic $O(n)$ loop model is defined by the generating function 
\begin{align}
    F^{(n)}_0[\mu] &= \sum_{k=3}^\infty\frac{1}{k!} \int_{0}^\infty \rmd \ell_1 \mu(\ell_1) \cdots \int_{0}^\infty \rmd \ell_k \mu(\ell_k)
	\sum_\Gamma \int_{\mathcal{M}_{0,k}(\ell_1,...,\ell_k)}\rmd_{WP} \prod_{\gamma\in \Gamma} n \ e^{-x L_\gamma}
\end{align}
where $n\in(0,2)$, $\Gamma$ is a collection of non-intersecting simple geodesic loops $\gamma$ of length $L_\gamma$ and $\rmd_{WP}$ is the Weil-Petersson measure of the hyperbolic surface. The coupling constant $n$ introduces a weight $n\ e^{-x L_\gamma}$ to each loop drawn on top of a hyperbolic surface (see Figure \ref{fig:hyper_o(n)}).\\  
What we aim to investigate is if there is a consistent weight substitution $\mu \rightarrow\hat{\mu}$ that relates the disk functions of the hyperbolic $O(n)$ loop model and the hyperbolic non-generic critical surfaces one, $W_\mu^{(n)}(\ell)=W_{\hat{\mu}}(\ell)$. From \eqref{eq:disk_non-gen-crit}, we can observe that asymptotically
\begin{equation}
    \hat{\mu}(\ell)= 2\;\sin\left(\pi\left(\alpha-\frac{1}{2}\right)\right)  \ e^{-2\sqrt{E_\mu}\ell} \ell\; 2^{-2}\; W_{\hat{\mu}}(\ell)  \label{weigth_substitution_hypo(n)}
\end{equation}
Therefore, we have shown there is a weight substitution that  
\begin{equation}
    \hat{\mu}(L_\gamma) = \mu(L_\gamma) + n  \ e^{-L_\gamma x} L_\gamma \; 2^{-2}\;W_{\hat{\mu}}(L_\gamma)\label{eq:on_non_gen_critical}
\end{equation}
for $x=2\sqrt{E_\mu}$ and $n=2\sin\left(\pi\left(\alpha-\frac{1}{2}\right)\right)$\footnote{The factor of $2^{-2}$ is necessary for consistency with the prefactor in the generating function of Weil-Petersson volumes and in \eqref{eq:JT_disk_function}. Thus, it does not affect the relation between $n$ and $\alpha$.}. In geometrical terms, this means that one can substitute each part of the hyperbolic surface surrounded by a geodesic loop $\gamma$ by a geodesic boundary with weight \eqref{eq:on_non_gen_critical} (see Figure \ref{fig:holes_maps}).\\
\begin{figure}[htbp]
\centering
\includegraphics[width=\textwidth]{hyper_on_full.pdf}
\caption{\textbf{Left}: Example of a hyperbolic $O(n)$ loop model with two non-intersecting geodesic loops in red. \textbf{Right}: The corresponding gasket where the marked geodesic boundaries are colored in gray. \label{fig:hyper_o(n)}}
\end{figure}\\
To conclude this chapter, we present the qualitative phase diagram of the $O(n)$ loop model (Figure \ref{fig:phase_diagram}) as studied in \cite{Borot:2011kg, Borot_2012}, which we expect to generalise for the hyperbolic $O(n)$ loop model. This will be relevant when we identify non-generic criticality in JT gravity (Section \ref{Section_Criticality}) and draw  similarities with hyperbolic geometries. \\
As we can see in Figure, \ref{fig:phase_diagram}, within the non-generic critical family of weights, there are two cases, the \textbf{\textit{dense}} and \textbf{\textit{dilute}} phases. In the dense phase ($\alpha>2$), the length of the loops is big and they can touch tangentially. In the dilute phase ($\alpha<2$), the loops tend to avoid each other and also themselves. 
\begin{figure}[htbp]
\centering
\includegraphics[width=.8\textwidth]{phase_diagram_on.pdf}
\caption{Qualitative phase diagram of the $O(n)$ model \cite{Borot:2011kg,Borot_2012}. The line $x=2\sqrt{E_\mu}$ corresponds to non-generic criticality, $t>0$ and $\mu_0(\ell)$ is a reference weight. \label{fig:phase_diagram}}
\end{figure}
%%%%%%%%%%%%%%%%%%%%%%%%%%%%%%%%
\section{JT Gravity}\label{section_JT_random}
Let us turn now to the standard formulation of JT gravity before reintroducing it in terms of random geometry. JT gravity is a 2-dimensional theory of gravity defined by the action 
\begin{equation}
    I_{JT}=-\frac{S_0}{2\pi}\left(\int_\mathcal{M}\sqrt{g}R+\int_{\partial\mathcal{M}}\sqrt{h}K\right)-\left(\frac{1}{2}\int_\mathcal{M}\sqrt{g} \phi (R+2) +\int_{\partial\mathcal{M}} \sqrt{h} \phi (K-1)\right)\label{eq:JTaction},
\end{equation}
where the first term is the Einstein-Hilbert action, which is topological in 2D and is proportional to the Euler characteristic of the surface $\mathcal{M}$, $\chi(\mathcal{M})=2-2g-n$, where $g$ is the genus of the surface and $n$ is the number of boundaries. The second term contains a scalar field $\phi$, called \textit{dilaton}, which acts as a Lagrange multiplier that imposes a negative curvature $R=-2$ in the bulk, this is called AdS spacetime. This is equivalent to imposing $\mathcal{M}$ to be a hyperbolic surface of genus $g$ and a boundary $\partial \mathcal{M}$ composed by a disjoint union of $n$ asymptotic boundaries with extrinsic curvature $K$ \cite{Saad:2019lba}. It is worth noting that the term ``asymptotic boundary" accounts for the fact that AdS spacetime is infinite in extent and does not have a traditional boundary that restricts or limits the spacetime itself but has a conformal boundary located at infinity and serves as a reference for defining the conformal structure of the spacetime.\\
Hence, the $n$-boundary JT gravity partition function reduces to a sum over geometries with asymptotic boundaries of length $(\beta_1,\dots, \beta_n)$ which, after integrating out the dilaton in the bulk, is
\begin{equation}
    Z(\beta_1,\dots,\beta_n)=\int_\mathcal{M} \mathcal{D}g\hspace{1mm}\mathcal{D}\phi \hspace{1mm} e^{-I_{JT}}= \int_{\mathcal{M}}
  \mathcal{D}g \hspace{1mm} e^{-S_0 \chi(\mathcal{M})} \delta(R+2) \int\mathcal{D}\phi \hspace{1mm} e^{\int_{\partial\mathcal{M}} \sqrt{h} \phi (K-1)}.\label{eq:JTpathintegral}
\end{equation}
The last expression hints at a topological expansion of spacetimes that split between bulk and boundary dynamics. It geometrically translates to the existence of a geodesic closed curve of length $b_i$ per asymptotic boundary $\beta_i$ that split a hyperbolic surface into two parts: bulk and \textit{trumpets}. The bulk part of \eqref{eq:JTpathintegral} is an integral over the moduli space of hyperbolic surfaces of genus $g$ and $n$ geodesic boundaries, and it is given by the Weil-Petersson volume $V^{WP}_{g,n}(b_1,\dots,b_n)$ \cite{Mirzakhani}. 
The trumpet is the part of the geometry that connects the geodesic boundaries with the asymptotic boundaries. The asymptotic boundary term in \eqref{eq:JTpathintegral} is given by the Schwarzian derivative of the near-boundary coordinates, which encodes information about the character of the asymptotic conformal symmetry group. Additionally, from the holographic point of view, it effectively describes a one-dimensional quantum mechanical system given by the so-called Schwarzian action. Then, the trumpet partition function is the partition function for a geometry with the Schwarzian action in the asymptotic boundary ending on a geodesic boundary, and it is given by
\begin{equation}
    Z_{tr}(\beta,b)=\frac{e^{-b^2/4\beta}}{2\sqrt{\pi \beta}}.
\end{equation}
Therefore, \eqref{eq:JTpathintegral} can be written as 
\begin{equation}
    Z(\beta_1,\dots,\beta_n)=\sum_{g=0}^\infty e^{-S_0 \chi}\int_0^\infty b_1 \rmd b_1 \dots \int_0^\infty b_n\rmd b_n V_{g,n}(b_1,\dots,b_n)Z_{tr}(\beta_1,b_1)\dots Z_{tr}(\beta_n,b_n).\label{eq:JTpart_g_n}
\end{equation}
From now on we will focus on the case $n=1$, i.e.\ one asymptotic boundary. In this case, \eqref{eq:JTpart_g_n} becomes 
\begin{equation}
    Z(\beta)=\sum_{g=0}^\infty e^{-S_0(2g-1)} \tilde{Z}_g(\beta)= e^{S_0} \tilde{Z}_0(\beta) + \sum_{g=1}^\infty e^{-S_0(2g-1)}\int_0^\infty b\rmd b V^{WP}_{g,1}(b)Z_{tr}(\beta,b).
\end{equation}
We separate the contribution for $g=0$ because $V_{0,1}$ vanishes. However, this geometry corresponds to a hyperbolic disk for which the path integral \eqref{eq:JTpathintegral} is not vanishing, this is what we call $\tilde{Z}_0(\beta)$ in this case, and it is given by \cite{Saad:2019lba}
\begin{equation}
    \tilde{Z}_0(\beta)=\frac{1}{\sqrt{16\pi\beta^3}}e^{\frac{\pi^2}{\beta}}
\end{equation}
%%%%%%%%%%%%%%%%%%%%%%%%%%%%%%%%%%%%%%%%
\subsection{JT Gravity with defects}
In a similar fashion, one can consider the case when the hyperbolic surface has geodesic boundaries of length $\ell$ that are not attached to an asymptotic boundary by a trumpet, these are called \textit{defects} \cite{Turi:2020ale} and are assigned an effective weight $\mu(\ell)$ that compile the physical system they represent e.g. conical defects or branes (See Fig. \ref{fig:branes}). 
\begin{figure}[htbp]
\centering
\includegraphics[width=.9\textwidth]{configuration_brane.pdf}
\caption{Left: JT spacetime ending on $n$ branes. Right: JT spacetime with $n$ geodesic boundaries with weight $\mu$. \label{fig:branes}}
\end{figure} 
Then, the partition function can be expanded in a sum over genus and a sum over the number of defects 
\begin{align}
    Z(\beta) =\sum_{g=0}^\infty e^{-S_0(2g-1)}\sum_{k=0}^{\infty} e^{-S_0 k} Z_{g,k+1}(\beta).\label{eq:full_JT_partfunction}
\end{align}
where for $g>0$ and $k\geq 0$,
\begin{equation}
    Z_{g,k+1}(\beta)=\frac{1}{k!}\int_0^\infty \rmd \ell_1
\mu(\ell_1)...\int_0^\infty \rmd \ell_k
\mu(\ell_k)\int_0^\infty b\rmd b V^{WP}_{0,k+1}(b,\ell_1,...,\ell_k)Z_{tr} (\beta,b)
\end{equation}
In particular, for $g=0$ we know that the Weil-Petersson volumes vanish for the cases $k=0,1$, but these partition functions can be explicitly computed using the Schwarzian action and they correspond to the disk ($\tilde{Z}_0$) and the trumpet ending on a defect (sometimes called \textit{half-wormhole}). Therefore, the genus-$0$ partition function can be explicitly written as 
\begin{align}
    Z_0(\beta)
&=\frac{1}{\sqrt{16\pi\beta^3}}e^{\frac{\pi^2}{\beta}}+e^{-S_0}\int_0^\infty \rmd b\mu(b)\frac{e^{-b^2/4\beta}}{2\sqrt{\pi \beta}}\\ 
    &+ \sum_{k=2}^\infty \frac{e^{-S_0k}}{k!}\int_0^\infty \rmd \ell_1
\mu(\ell_1)...\int_0^\infty \rmd \ell_k
\mu(\ell_k)\int_0^\infty b\rmd b\frac{e^{-b^2/4\beta}}{2\sqrt{\pi \beta}}V^{WP}_{0,k+1}(b,\ell_1,...,\ell_k).\label{eq:full_diskamplitude}
\end{align}
More examples of the geometries that contribute to \eqref{eq:full_JT_partfunction} are shown in Figure \ref{fig:JT_full_sum}.\\
\begin{figure}[htbp]
\centering
\includegraphics[width=\textwidth]{configuration.pdf}
\caption{Examples of geometries that contribute to the partition function \eqref{eq:full_JT_partfunction}. From left to right: $(g=0,n=1,k=0)$, $(g=0,n=1,k=1)$, $(g=2,n=1,k=0)$ and $(g=2,n=1,k>0)$. \label{fig:JT_full_sum}}
\end{figure}\\
So far, we have introduced JT gravity as it is mainly done in the physics literature, yet, one of our main goals is to reintroduce techniques from random geometry that we will prove are highly useful for JT gravity (Section \ref{Section_Criticality}). In particular, the power of the string equation formulation has been used in JT gravity \cite{Turiaci:2020fjj,Forste:2021roo}\footnote{It is worth noting that what we call the string equation is known in the physics literature as the genus $g=0$ string equation in the limit $\hbar\rightarrow 0$}, but, in the physics literature, one usually computes the partition function and density of states using a holographic quantum mechanical system in the asymptotic boundary and then `guess' the form of the string equation. Here, we take the opposite point of view and start with the string equation.\\
% \subsection{Generating function and string equation}
% So far, we have introduced JT Gravity as it is mainly done in the Physics literature, yet, one of our main goals is to reintroduce techniques from random geometry that we prove are highly useful for JT Gravity in Section \ref{Section_Criticality}. Following the same reasoning we presented in Section \ref{section:Random_geometry}, we want to consistently define a probability measure in the space of hyperbolic surfaces of genus $g$ and $n$ geodesic boundaries, \Anote{NOTATION?} $\mathcal{M}^{WP}_{g,n}$. This is done by using the Weil-Petersson volumes in the same way as for Boltzmann genus-$g$ maps. The generating function of Weil-Petersson volumes of genus $g$ is defined as \cite{1996CMaPh.181..763K}
% \begin{align}
% 	F^{WP}_g[\mu] &= \sum_{k=3}^\infty\frac{2^{2-2g-k}}{k!} \int_{0}^\infty \rmd \alpha_1 \mu(\alpha_1) \cdots \int_{0}^\infty \rmd \alpha_k \mu(\alpha_k) V^{WP}_{g,k}(\alpha_{1}\ldots,\alpha_{k}),
% \end{align}
% where we set $e^{S_0}=1$ for convenience. It can be expressed in terms of the string equation
% \begin{equation}
%     x(u)=\frac{\sqrt{u}}{\pi}J_1\left(2\pi\sqrt{u}\right)-\int_0^\infty \rmd\ell \mu(\ell) I_0\left(\ell \sqrt{u}\right)\label{eq:general_string_eq_JT}
% \end{equation}
% which is well-known in the physics literature up to a change of sign $u\rightarrow -u$. Equivalently, it can be written in the form 
% \begin{equation}
%    x(u)=u-\sum_{m=0}^\infty\frac{t_m+\gamma_m}{m!}u^m,
% \end{equation}
% with 
% \begin{equation}
% \gamma_k=\frac{(-1)^k\pi^{2k-2}}{(k-1)!}\mathrm{1}_{k\geq 2}\label{gamma_k}
% \end{equation}
% the KdV times for pure JT Gravity [original ref?]\cite{Johnson:2019eik,Okuyama:2019xbv} and
% \begin{equation}
%     t_k=\frac{4^{-k}}{k!}\int_0^\infty \rmd \ell \mu(\ell) \ell^{2k}\label{eq:def_general_JT_times}
% \end{equation}
% the times associated to geodesic boundaries with weights $\mu(\ell)$. Note that for Boltzmann metric maps, we have $\gamma_k=0$.\\
% The power of the string equation formulation has been used in JT Gravity \cite{Turiaci:2020fjj,Forste:2021roo}\footnote{It is worth noting that what we call the string equation is known in the physics literature as the genus $g=0$ string equation in the limit $\hbar\rightarrow 0$}, but, in the physics literature, one usually computes the partition function and density of states using a holographic quantum mechanical system in the asymptotic boundary and then `guess' the form of the string equation. Here, we take the opposite point of view and start with the string equation.\\
% From now on, we will restrict our study to the case $g=0$ for convenience, but the interested reader can check \cite{budd2022irreducible,budd2023_2} for $g>0$. The genus $g=0$ generating function of Weil-Petersson volumes can be expressed in the same way as for Boltzmann planar metric maps, as well as the moment $M_0$
% \begin{align}
%     	F^{WP}_0[\mu] &= \frac{1}{4}\int_{0}^{E_\mu} \!\!\!\rmd u\, x(u)^2,\label{eq:genfun0}\\
%     M^{WP}_0&=\frac{\partial x}{\partial u} (E_\mu)\label{eq:m0_wp}
% \end{align}
% where $E_\mu>0$ is the smallest solution to $x(E_\mu)=0$, and, in the physics literature, it is called the threshold energy $E_\mu=-E_0$.  This allows us to define a probability density measure on the space of hyperbolic genus-$0$ surfaces  
% \begin{equation}
%     p^{WP}\left(\ell_1, ... , \ell_n\right) = \frac{1}{n!} \frac{\prod_{i=1}^{n}  \mu(\ell_i)}{F^{WP}_0[\mu]}  V^{WP}_{0,n}(\ell_1,\dots,\ell_n).
% \label{eq:prob_density_WP}
% \end{equation}
% In this way, we can call \textbf{\textit{random genus-$0$ hyperbolic surface with $n$-boundaries}} to a surface sampled with probability \eqref{eq:prob_density_WP}.\\
% In particular, it is useful to introduce the disk function, this is the generating function of hyperbolic surfaces with one marked boundary of length $b$
% \begin{equation}
% W(b)=\frac{\delta F
% _0}{\delta\mu(b)}=\sum_{k=2}^\infty\frac{2^{2-(k+1)}}{k!}\int_0^\infty \rmd \ell_1
% \mu(\ell_1)...\int_0^\infty \rmd \ell_k
% \mu(\ell_k)V^{WP}_{0,k+1}(b,\ell_1,...,\ell_k).
% \end{equation}
Using this, we can re-express \eqref{eq:full_diskamplitude} in a more convenient way which explicitly makes use of the string equation 
% \begin{align}
%     Z_0(\beta)
% &= \frac{1}{\sqrt{16\pi\beta^3}}e^{\frac{\pi^2}{\beta}}+\int_0^\infty b\rmd b\frac{e^{-b^2/4\beta}}{2\sqrt{\pi \beta}}\frac{\mu(b)}{b}+\int_0^\infty b\rmd b\frac{e^{-b^2/4\beta}}{2\sqrt{\pi \beta}}W(b).\label{disk}
% \end{align}
\cite{Johnson:2019eik,Okuyama:2019xbv,Turi:2020ale}
\begin{align}
 Z(\beta)_{g=0}&=\frac{1}{2\sqrt{\pi\beta}}\int_{-\infty}^0\rmd x e^{u(x) \beta }=
\frac{1}{2\sqrt{\pi\beta}}\int_{-\infty}^{E_\mu}\rmd u e^{u \beta }\frac{\partial x(u)}{\partial u}\label{eq:Z0def}\\
&=  -\frac{1}{2}\int_{-\infty}^0\rmd u e^{u \beta }x(u)\sqrt{\frac{\beta}{\pi}} -\frac{1}{2}\int_{0}^{E_\mu}\rmd u e^{u \beta }x(u)\sqrt{\frac{\beta}{\pi}}\\
&=  \frac{1}{\sqrt{16\pi\beta^3}}e^{\frac{\pi^2}{\beta}}+\int_0^\infty \rmd b\mu(b)\frac{e^{-b^2/4\beta}}{2\sqrt{\pi \beta}} + \int_0^\infty b\rmd b\frac{e^{-b^2/4\beta}}{2\sqrt{\pi \beta}}W(b),
\end{align}
where we set $e^{S_0}=1$ for convenience and in the last line we used the identity
\begin{equation}
    e^{-bz}=2z\int_0^\infty \rmd \beta e^{-\beta z^2}\frac{e^{-b^2/4\beta}}{2\sqrt{\pi \beta}}.
\end{equation}
The density of states is obtained by Laplace transforming the partition function, and is given by
\begin{equation}
        \rho_{0}(E) = \frac{1}{2\pi}\int_{-E}^{-E_0}\rmd u \frac{1}{\sqrt{E+u}} \frac{\partial x}{\partial u}\label{eq:rho0def}.
\end{equation}
In order to draw a parallel between this language and some well-known results in the JT gravity literature, we present two examples: FZZT and end-of-the-world branes.
\subsection{Example: FZZT Branes}
Dynamical FZZT Branes were introduced in the context of Liouville Field Theory in \cite{Fateev:2000ik} as a family of conformally invariant (Dirichlet) boundary conditions parametrized by the so-called \textit{boundary cosmological constant}, which controls the length of the boundary. Beyond the Physics aspects, the convenience of working with these branes comes from the fact that they have well-defined analytic expressions in Matrix Models, they are determinant insertions \cite{Aganagic:2003qj}. In the JT gravity literature, FZZT branes have been studied in this language \cite{Saad:2019lba,Blommaert:2019wfy}, and it was shown that these branes can be effectively encoded by geodesic boundaries of length $\ell$ with weight \cite{Okuyama:2021FZZT} (see Figure \ref{fig:branes})
\begin{equation}
    \mu_{FZZT}(\ell)=-e^{-z \ell}.
\end{equation}
where the constant $z$ plays the role of the boundary cosmological constant. Using \eqref{eq:def_general_JT_times}, we can compute the times 
\begin{equation}
    t_k^{FZZT}=-\frac{4^{-k}(2k)!}{k!}z^{-2 k-1}, 
\end{equation}
which are consistent with \cite{Okuyama:2021FZZT}. Then, the string equation \eqref{eq:general_string_eq_JT} takes the well-known form
\begin{equation}
     x(u)=\frac{\sqrt{u}}{\pi}J_1\left(2\pi\sqrt{u}\right)+\frac{1}{\sqrt{z^2-u}}.
\end{equation}
It is worth noting that the term `dynamical' in this context means that all physical quantities, e.g. \eqref{eq:Z0def} and \eqref{eq:rho0def}, are computed using the shifted weights $t_m+\gamma_m$ instead of the pure JT gravity ones, $\gamma_m$. In these terms is that our computations are also done in the dynamical framework.\\
% and dilaton potential 
% \begin{equation}
%     U(\Phi)=2\Phi -\frac{z}{\Phi ^2+z^2}e^{-2\pi\Phi}
% \end{equation}
% \subsection{Example: Conical Defects}
% This has been generalized to JT Gravity with conical defects, where the conical defects are equivalent to non-asymptotic boundaries of imaginary length, in \cite{Turi:2020ale} it is defined 
% \begin{equation}
%     \mu(\ell)=\lambda \delta(\ell-i 2\pi\alpha)
% \end{equation}
\subsection{Example: EOW branes}
In JT gravity, End-of-the-world branes are boundaries of space where Neumann boundary conditions are usually imposed. These have been mainly used to investigate black hole evaporation and the Page curve\cite{Penington:2019kki}, and correspond to introducing vectors in the JT gravity Matrix Model. However, for our purposes, EOW branes are geodesic boundaries with weights \cite{Gao:2021uro}
\begin{equation}
    \mu_
{EOW}(\ell)=\frac{e^{-m\ell}}{2\sinh{\left(\ell/2\right)}}
\end{equation}
and, equivalently, times 
\begin{equation}
    t^{EOW}_k=\frac{4^{-k}(2k)! \zeta \left(2 k+1,m+\frac{1}{2}\right)}{k!} \mathrm{1}_{k\geq 1},
\end{equation}
where $m$ is the value of the normal derivative of the dilaton at the EOW boundary.
% \begin{align}
%     Z(\beta)_{g=0} &=\frac{1}{\sqrt{16\pi\beta^3}}e^{\frac{\pi^2}{\beta}}+\int_0^\infty \rmd b\mu(b)\frac{e^{-b^2/4\beta}}{2\sqrt{\pi \beta}}\\ 
%     &+ \sum_{k=0}^\infty  \frac{\lambda^k}{k!}\int_0^\infty b \rmd b \frac{e^{-b^2/4\beta}}{2\sqrt{\pi \beta}} V_{0,1,k}(b,2\pi i a_1,...,2\pi i a_k)
% \end{align}

% \begin{equation}
%     W(\Phi)=2\Phi+\lambda e^{-2\pi(1-\alpha)\Phi}
% \end{equation}

\section{Critical JT Gravity}\label{Section_Criticality}
Our objective is to investigate the applicability of the criticality results observed in Boltzmann planar maps and more recently in hyperbolic surfaces \cite{budd2023} to JT gravity. In this last work, it is proven that the generating function of Weil-Petersson volumes experiences a similar critical behavior to Boltzmann planar maps (see Table \ref{tab:critical}). Furthermore, in the non-generic critical regime ($M_0^{WP}=0$), the hyperbolic surfaces also develop large faces, which correspond to geodesic boundaries. We are specially interested in how the non-generic criticality in the bulk translates into the holographic quantum system on the asymptotic boundary, and, in particular, how it affects its partition function and density of states. More precisely, we will show that this family of non-generic critical models consistently interpolates between two known cases in the JT literature, $\rho(E)\sim \sqrt{E}$ and $\rho(E)\sim E^{3/2}$, which correspond to the $k=1$ and $k=2$ minimal models respectively.\\
For simplicity, we will assume a general form of the non-generic critical weight $\mu_{\alpha}$ that satisfies the conditions deduced in Section \ref{section:Random_geometry} such that
% \begin{equation}
%     \mu_{ngc}(\ell)=-
%     \left(\frac{\delta (\ell+2 \pi i)}{\ell}+\frac{\delta (\ell-2 \pi i)}{\ell}\right)' +\mu_{\alpha}(\ell)\label{muform},
% \end{equation}
the string equation is  
\begin{equation}
    x_{crit}(u)= -t_{crit}\hspace{1mm}(E_\mu-u)^{\alpha -\frac{1}{2}} + \mathcal{O}((E_\mu-u)^{\alpha -\frac{1}{2}})\label{muform}
\end{equation}
where $3/2<\alpha<5/2$. 
The genus-$0$ generating function is
\begin{equation}
    F^{WP}_0[\mu_\alpha]=\frac{E_\mu^{2 \alpha }}{8 \alpha } t_{crit} +\mathcal{O}(E_\mu^{2\alpha}),
\end{equation}
and the disk function 
\begin{equation}
    W(b)=\frac{E_\mu^{\alpha +\frac{1}{2}} \, _0F_1\left(\alpha +\frac{3}{2};\frac{b^2 E_\mu}{4}\right)}{2 \alpha +1} t_{crit} + \mathcal{O}(b^{-\alpha-1}).
\end{equation}
% while the moments are 
% $M_{p-1}=(-1)^{p+1} \left(\alpha -\frac{1}{2}\right)^{(p)} (E_\mu-u)^{\alpha -p-\frac{1}{2}}$.
\begin{figure}[htbp]
\centering
\includegraphics[width=0.8\textwidth]{configurationcopy.pdf}
\caption{Geometries that contribute to the partition function of critical JT gravity with defects of weight $\mu_\alpha$.\label{fig:criticalJT}}
\end{figure}\\

\subsection{Density of states}
As emphasized in Section \ref{section_JT_random}, the string equation provides sufficient information for determining the leading-order genus zero partition function and density of states. So, using \eqref{muform}, \eqref{eq:Z0def} and \eqref{eq:rho0def}, and neglecting terms of higher order, the partition function of critical JT gravity is 
\begin{equation}
    Z_0(\beta)= \frac{ \Gamma \left(\alpha +\frac{1}{2}\right) e^{-\beta  E_0}t_{crit}}{2 \sqrt{\pi }}\beta ^{-\alpha }\label{eq:Z_beta},
\end{equation}
and its density of states is given by \eqref{eq:rho0def}, 
\begin{equation}
     \rho_{0}(E) = \frac{\Gamma \left(\alpha +\frac{1}{2}\right) t_{crit}}{2 \sqrt{\pi } \Gamma (\alpha )}(E-E_0)^{\alpha -1}\label{eq:rho_E},
\end{equation}
where we substituted $E_\mu\rightarrow E_0$. In Figure \ref{fig:critical_density}, we show \eqref{eq:rho_E} for different values of $\alpha$. We note that this non-generic critical family of models interpolates constantly between $\sqrt{E-E_0}$ and $(E-E_0)^{3/2}$.\\
\begin{figure}[htbp]
\centering
\includegraphics[width=.8\textwidth]{density_energy.pdf}
\caption{Genus zero density of states \eqref{eq:rho_E} for different values of $\alpha$ and $t_{crit}=1$. This family of non-generic critical models continuously interpolates between $\sqrt{E-E_0}$ and $(E-E_0)^{3/2}$. Our simplified forms of the weight and string equation do not fix a value for $E_0$. For convenience, we consider the threshold energy as constant ($E_0=0.1$) for all $\alpha$, but, in general, it will also be a function of $\alpha$.  \label{fig:critical_density}}
\end{figure}

\subsection{Phase transitions}
We dedicate this subsection to clarify the physical and geometrical meaning of the JT gravity non-generic critical family of models described by the density of states \eqref{eq:rho_E}. First, let us remark that it is known that in JT gravity with FZZT \cite{Okuyama:2021Page} and EOW \cite{Gao:2021uro} branes  some parameters can be tuned to shift $\rho_0(E)$ from $\sqrt{E-E_0}$ to $(E-E_0)^{3/2}$, and our model bridges the gap between these two cases by providing an interpolation. However, the framework that we present in this paper makes clear and precise in which sense there is a phase transition. The explanation is that by fine-tuning the first derivative of the string equation to vanish, the weight goes from a subcritical to a generic critical phase. What our model does is to fine-tune the weight to be non-generic critical. In principle, this non-generic critical weight can be linked to modifications to the FZZT and EOW brane weights. Given that our model offers this interpolation, it seems plausible this non-generic critical weight can be linked to deformations of the FZZT and EOW brane weights and its precise form could yield interesting results. However, we will not explore this direction in the present work, but rather the general properties of these deformations.\\
The aspect we highlight now is the geometrical features of this family of non-generic critical models. Within this critical regime, the system undergoes significant transformations. Notably, the analysis in \cite{budd2023} reveals formal divergences in the expected number of boundaries and their corresponding total length, indicating substantial geometrical similarities with the Boltzmann planar map case as the model approaches criticality.
% no, but one would expect so because you pass from loops that do not touch to ones that do. It is also expected to be SLE(4?). There is a correspondence between $\alpha$, n and $\gamma$ and $\kappa$\\
% write what is string equation from the planar maps case. Tree bijection? 
We can see this by using the probability measure \eqref{eq:prob_density_WP} to compute the expected number of geodesic boundaries  
        \begin{equation}
           \mathbb{E}(n) = \sum_{n=3}^{\infty} n\cdot\frac{2^{2-n}}{n!} \frac{\prod_{i=1}^{n} \int_{0}^{\infty}\rmd \ell_i \mu(\ell_i) V^{WP}_{0,n}(\ell_1 , ... , \ell_n)}{F^{WP}_0},
        \end{equation}
%while the expected total length of the geodesics boundaries $\ell_{\text{total}}$ is 
%\begin{equation}
 %   \begin{split}
  %      \mathbb{E}(\ell_{\text{total}}) &= \sum_{n=3}^{\infty} p_n(\ell_1, ... , \ell_n)\cdot \left(\ell_1 + ... + \ell_n\right) \\
   %     &= \sum_{n=3}^{\infty} \frac{2^{2-n}}{n!} \frac{\prod_{i=1}^{n} \int_{0}^{\infty}\rmd \ell_i \ \mu(\ell_i)}{F^{WP}_0} V^{WP}_{0,n}(\ell_1 , ... , \ell_n) \ \left(\ell_1 + ... + \ell_n\right) 
    %\end{split}
    %\label{eq: E(L_tot)}
%\end{equation}
The main argument is that when two of the defects have length zero, we can easily see that  
\begin{equation}
    \begin{split}
        \frac{\delta E_\mu}{\delta \mu(\ell)} =
        \sum_{n=3}^{\infty} \frac{2^{3-n}}{(n-3)!} \left( \prod_{k=4}^{n} \int_{0}^{\infty}  \rmd \ell_i \ \mu\left(\ell_i\right) \right) V^{WP}_{0,n}\left(0, 0, \ell, \ell_4, ..., \ell_{n}\right) \rightarrow \infty
    \end{split}
    \label{eq: R dmu}
\end{equation}
since, from the string equation, we have that $ \frac{\delta E_\mu}{\delta \mu(\ell)} = \frac{\delta u/\delta \mu (E_\mu)}{M^{WP}_0}$, and $M^{WP}_0=0$ at criticality. Therefore, $\mathbb{E}(n)\rightarrow\infty$.
Additionally, we can analyse the expected number of boundaries in a disk of diameter $b$ in the non-generic critical regime,
\begin{equation}
    \mathbb{E}_{W_\mu(b)}(n)=\int_0^\infty \rmd \ell \frac{\delta W_\mu(b)}{\delta \mu (\ell)}\mu(\ell) 	\propto \int_0^\infty \rmd \ell \; \ell^{-1/2}e^{\ell\sqrt{E_\mu}}\; \ell^{-\alpha}e^{-\ell \sqrt{E_\mu}} + \mathcal{O}(\ell^{-\alpha-1/2}).
\end{equation}
The expression inside the integral is the large $\ell$ probability distribution of the number of geodesic boundaries of size $\ell$ in a disk of diameter $b$. Then, we can see that the variance of this distribution is $\text{Var}(\ell^{-\alpha-1/2})\rightarrow\infty$ for $\alpha>3/2$. This implies that in the non-generic critical regime of JT gravity, one can expect a very large number of defects, and if $b\rightarrow\infty$ there will always be boundaries of comparable size. 
Even though we are not explicitly proving the existence of a dilute and dense phase, we anticipate a similar behavior to Boltzmann maps due to the relationship with the hyperbolic $O(n)$ loop model discussed in Section \ref{section:Random_geometry}. A schematic illustration of this behavior is provided in Figure \ref{fig:holes_JT}.\\
  \begin{figure}[htbp]
   \centering
   \includegraphics[width=.9\textwidth]{holes_JT.pdf}
   \caption{Schematic illustration of a hyperbolic surface with: subcritical, dilute, and dense non-generic critical weights.\label{fig:holes_JT}}
  \end{figure}
% Some phase transitions have been identified within the current JT Gravity literature. These include the transition from a one-cut matrix model (MM) to a multi-cut MM, as well as the observed transitions in end-of-the-world (EOW) branes and FZZT branes from $E^{1/2}$ to $E^{3/2}$. These findings carry significance due to the suitability of FZZT branes for description in terms of Matrix Models and their role as boundary conditions in Liouville Quantum Gravity. However, it should be noted that these phase transitions are not yet fully formalized, and their non-perturbative stability remains an open question. Additionally, the compatibility of these transitions with alternative definitions of EOW branes, such as those proposed in \cite{Johnson:2021tnl}, warrants further investigation.

% -Boundary conditions are in AdS boundary


% This is consistent with the phase transition found studying end-of-the-world branes in \cite{Gao:2021uro}
% Both for FZZT and EOW branes there is a phase transition from $\sqrt{E-E_0}$ to $(E-E_0)^{3/2}$ and our model interpolates between these. 


% -FZZT Branes come from Liouville QG 
% and have been extended to JT Gravity

% And a behaviour for the density of states was also analysed in \cite{Okuyama:2021Page} when computing black hole entropy.


% -Idea: Is there a modification of the EOW brane weight that works? Sounds like it should be since this model interpolates between two EOW cases. \cite{Gao:2021uro}



\section{Critical Dilaton Gravity}\label{Section_Dilaton_gravity}
In the field of JT gravity with defects, such as conical defects and branes,  it has been proposed that these can be integrated out, leading to an effective dilaton action \cite{Witten:2020wvy, Turi:2020ale, Gao:2021uro}, which is a deformation of pure JT gravity. While the formality of this approach may remain subject to an ongoing debate, we adopt this particular conjectural view as a means to analyze non-generic critical models within the field of thermodynamic solutions in dilaton gravity.\\
Dilaton gravity is a 2-dimensional theory defined by the action 
\begin{equation}
    I_{dil}=\frac{1}{2}\int_{\mathcal{M}}\sqrt{g}(\Phi R+W(\Phi)).\label{eq:dilaton_action}
\end{equation}
where $\Phi$ is the dilaton and the potential is given by\footnote{\cite{Okuyama:2021FZZT} highlights the caveats of constructing the dilaton potential for a time modification different from the conical defect.}
    \begin{equation}    W(\Phi)=2\Phi+e^{-2\pi\Phi}\int_0^\infty \rmd \ell \mu(\ell)\cos\left(\ell \Phi\right).\label{eq:dilaton_pot_general}
    \end{equation}
The first term corresponds to the pure JT gravity case and the second term is determined by the weight per geodesic boundary. This potential satisfies the asymptotically AdS$_2$ condition, i.e.\ $W(\Phi)\sim 2\Phi$ as $\Phi\rightarrow +\infty$. Yet, its closed form  is only known for a few cases e.g. FZZT branes. In the case of critical JT gravity, we use a generalization of the form of the non-generic critical weight we derived in Section \ref{section:Random_geometry}, $\mu_{\alpha}(\ell)= e^{-\ell \sqrt{E_\mu}}\ell^{-\alpha}$ to compute 
% \begin{equation}
%     W(\Phi)=2\Phi-\frac{\Phi}{\pi}e^{-2\pi\Phi}\sinh\left(2\pi\Phi\right)+e^{-2\pi\Phi}\int_0^\infty \rmd \ell \mu_{\alpha}(\ell)\cos\left(\ell\Phi\right)\label{JT_dilaton_potential_critical_general}
% \end{equation}
the dilaton potential
\begin{align}
    W(\Phi) &=2\Phi
    +t_{crit} e^{-2\pi\Phi}\sum_{m=0}^\infty \frac{(-1)^m}{(2m)!}\Phi^{2m} \left(\sqrt{E_\mu}\right)^{\alpha-2m-1}\Gamma\left(2m+1-\alpha\right),
    % \\
    % &= t_{crit}\Gamma (1-\alpha ) \left(E_\mu+\Phi ^2\right)^{\frac{\alpha -1}{2}} \cos \left((\alpha -1) \tan ^{-1}\left(\frac{\Phi }{\sqrt{E_\mu}}\right)\right)
    \label{eq:dilaton_pot_critical_particular}
\end{align}
This potential is shown in Figure \ref{fig:dilaton_potential} for different values of $\alpha$ and in Figure \ref{fig:dilaton_potential_t} for different values of $t_{crit}$ for the dilute and dense phases. \\
It should be emphasized that the weight $\mu_{\alpha}(\ell)= e^{-\ell \sqrt{E_\mu}}\ell^{-\alpha}$ is not finite at $\ell=0$, so it is not an admissible weight in the probabilistic sense. Nevertheless, we employ this specific form to examine the overall features of the dilaton potential. Specifically, this becomes significant when performing the analytical continuation of the gamma function $\Gamma(2m+1-\alpha)$, which encompasses the non-analytic behavior of the weight $\mu_\alpha$. The physical consequences of this will be explored in the following subsection.
%%%%%%%%%%%%%%%%%%%%%%%%%%%%%%%%
\subsection{Thermodynamics}
Following the steps of \cite{Witten:2020ert}, we can write the classical euclidean black hole solution to the equations of motion of \eqref{eq:dilaton_action} in the gauge $\Phi(r)=r$,
      \begin{equation}
       ds^2=A(r)^2dt^2+\frac{1}{A(t)^2}dr^2, \hspace{12mm} A(r)=\int_{r_h}^r \rmd x W(x)
      \end{equation}
where $r=r_h$ is the position of the euclidean horizon. To have a well-defined exterior solution, we need this function to be positive $A(r)>0$ for $r>r_h$ and vanish at the horizon. Expanding around the horizon one gets $A(r)=W(\Phi_h)(r-r_h)+\mathcal{O}((r-r_h)^2)$, therefore the latter condition translates into $W(\Phi_h)>0$. Substituting this expansion in the metric and demanding smoothness at $r=r_h$, one obtains the euclidean temperature of the black hole 
\begin{equation}
    T(\Phi_h)=\frac{W(\Phi_h)}{4\pi}.
\end{equation}
% An example is 
% \begin{equation}
% \frac{1}{2} t \exp (-2 \pi  \phi ) e^{\sqrt{R}-i \phi } \left(e^{2 i \phi } E_{\alpha }\left(i \phi +\sqrt{R}\right)+E_{\alpha }\left(\sqrt{R}-i \phi \right)\right)-\frac{\phi  \exp (-2 \pi  \phi ) \sinh (2 \pi  \phi )}{\pi }+2 \phi
% \end{equation}
Additionally, it can be observed that the energy difference between two black holes with horizon radii $\Phi(r_1) = \Phi_1$ and $\Phi(r_2) = \Phi_2$ is
\begin{equation}
    \Delta E=\frac{1}{2}\int_{\Phi_1}^{\Phi_2}\rmd \Phi \ W(\Phi).
\end{equation}
Thermodynamically speaking, a system is stable if its heat capacity is positive, i.e.\ $\frac{dE}{dT}>0$, which implies that the stability of the euclidean black hole solution with a horizon at $\Phi_h$ requires that $W'(\Phi_h)>0$. In the case of our potential \eqref{eq:dilaton_pot_critical_particular}, it is monotonically increasing in the dilute phase ($3/2<\alpha<2$) and convex in the dense phase ($2<\alpha<5/2$), as shown in Figure \ref{fig:dilaton_potential}. Moreover, this family of non-generic critical models continuously interpolates between these two regimes. A qualitatively similar behavior has been observed in \cite{Okuyama:2021Page}, but in that case, one regime corresponds to FZZT branes and the other to anti-FZZT branes and the density of states changes abruptly from $\sqrt{E-E_0}$ to $(E-E_0)^{3/2}$. In our case, the change in the dilaton potential is directly linked to the varying scaling of $\rho_0$.\\
\begin{figure}[htbp]
\centering
\includegraphics[width=.9\textwidth]{dilaton_E.pdf}
\caption{Dilaton potential \eqref{eq:dilaton_pot_critical_particular} for different values of $\alpha$ and $t_{crit}=1$.    \label{fig:dilaton_potential}}
\end{figure}
\begin{figure}[htbp]
\centering
\includegraphics[width=.8\textwidth]{dilaton_t.pdf}
\caption{Dilaton potential \eqref{eq:dilaton_pot_critical_particular} for different values of $t_{crit}$. The top four graphs correspond to $\alpha=2.25$ ($(E-E_0)^{5/4}$)and the bottom four correspond to $\alpha=1.75$ ($(E-E_0)^{3/4}$).  \label{fig:dilaton_potential_t}}
\end{figure}\\
The key remark we make is that these two conditions of the potential, $W(\Phi_h)>0$ and $W'(\Phi_h)>0$, define the physical consistency requirements for any non-generic critical potential we can construct, which are additional constraints for the form of the weights $\mu_\alpha(\ell)$. So, if we start with a general potential \eqref{eq:dilaton_pot_general} without assuming the exact form of $\mu_\alpha$ and want to construct thermodynamically stable black hole solutions, it implies that
\begin{equation}
     0<\exp^{-2\pi \Phi_h}\int_0^\infty \rmd \ell \ \ell\mu_{\alpha}(\ell)\sin{(\ell\Phi_h)}< 2+4\pi\Phi_h
\end{equation}
needs to be satisfied. It is to be remarked that this inequality admits weights that diverge at $\ell=0$, even though these are not allowed if we want a probabilistic interpretation of non-generic critical JT gravity. Furthermore, in terms of the physical interpretation of this family of non-generic critical models, they can be interpreted as a family of thermodynamic stable brane solutions with finite entropy, to which evaporating black holes can decay to, but this remains to be proven.

% \subsection{Phase transitions}\label{subSection_phase_transition}
% \Anote{rewrite}
% The understanding of phase transitions in the context of JT gravity comprises several aspects. the study of end-of-the-world (EOW) branes reveals another significant phase transition. These branes undergo a transition from a scaling behavior of $E^{1/2}$ to $E^{3/2}$, where $E$ represents the energy. This transition has implications for the dynamics of the branes and their role in the overall behavior of JT gravity. Similarly, FZZT branes, which can be described effectively using matrix models and have a connection to Liouville Quantum Gravity, also exhibit a phase transition from $E^{1/2}$ to $E^{3/2}$. The existence of this transition provides a valuable link between FZZT branes and the matrix model formalism, highlighting their relevance in understanding the properties of JT gravity.\\
% It is important to note that while these phase transitions have been observed, their complete formalization is still a subject of ongoing investigation. Non-perturbative stability and compatibility with alternative definitions of EOW branes, as explored in \cite{Johnson:2021tnl} and similar studies, remain open questions. Further research and analysis are required to establish the precise nature and stability of these phase transitions and their alignment with different brane definitions.
% % The currently known phase transitions are
% % \begin{itemize}
% %     \item from one-cut to multi-cut MM
% %     \item EOW branes from $E^{1/2}$ to $E^{3/2}$
% %     \item FZZT branes from $E^{1/2}$ to $E^{3/2}$. This is useful since FZZT branes have a good description in terms of Matrix Models and as Boundary conditions in Liouville Quantum Gravity.
% % \end{itemize}

% % These are not completely formal and it is still an open question if these are non-perturbatively stable and if they are compatible with another definitions of eow branes such as \cite{Johnson:2021tnl}.

% % \subsection{Matrix models}

% % \subsection{Branes}
% In this work, we follow the steps of \cite{Gao:2021uro}, where the effective zero temperature entropy is defined as the leading behavior of $\rho_0(E)$ near $E=E_0$, in the non-generic critical case, it corresponds to the prefactor of $(E-E_0)^{\alpha-1}$

% \begin{equation}
% e^{S_{0,\text{eff}}}=\frac{\Gamma\left(\alpha+\frac{1}{2}\right))}{2\sqrt{\pi\Gamma(\alpha)}} t_{crit}
% \end{equation}


% Furthermore,
% The understanding of phase transitions in the context of JT gravity comprises several aspects. Firstly, there is a phase transition observed from one-cut to multi-cut matrix models. This transition [description]

\section{Spectral form factor}\label{section:SFF}
The spectral form factor in JT gravity is given by the time-ordered two-point correlation function, which in geometrical terms corresponds to the double trumpet for genus zero and no defects. 
In matrix models, it displays distinctive and universal characteristics, notably the presence of a ramp followed by a plateau structure. Interestingly, in JT gravity the spectral form factor exhibits a ramp-like behavior without a well-defined plateau.
The absence of a plateau is a consequence of the geometric structure of the double trumpet, whose dominant contribution leads to a continuous growth of the spectral form factor without reaching a plateau. The absence of a plateau in the spectral form factor indicates that the dual CFT does not possess a well-defined time scale for thermalization.\\
In earlier studies (e.g. \cite{Okuyama:2021FZZT}), it has been investigated if a plateau can be recovered with the inclusion of branes. Hence, it arises as a natural question to investigate it using our configuration of branes.
In our case, the main contribution to the double trumpet corresponds to two trumpets (or half-wormholes) glued at geodesic boundaries with non-generic critical weight $\mu_{\alpha}(\ell)= e^{-\ell \sqrt{E_\mu}}\ell^{-\alpha}$,
\begin{align}
    Z_{0,2}(\beta_1,\beta_2) &= \int_0^\infty \rmd b\mu_\alpha(b)^2\frac{e^{-b^2/4\beta_1}}{2\sqrt{\pi \beta_1}} \frac{e^{-b^2/4\beta_2}}{2\sqrt{\pi \beta_2}}\\
    &=   \frac{t_{crit}}{4\pi}\; \frac{\Gamma (1-\alpha ) \left(\frac{\beta_1\beta_2}{\beta_1+\beta_2}\right)^{\frac{1-\alpha }{2}} U\left(\frac{1}{2}-\frac{\alpha }{2},\frac{1}{2},E_\mu\frac{ \beta_1 \beta_2}{\beta_1+\beta_2}\right)}{ \sqrt{\beta_1\beta_2}},
\end{align}
where $U$ is the confluent hypergeometric function. Therefore, the spectral form factor is 
\begin{equation}
    Z_{0,2}(\beta+it,\beta-it)=\frac{t_{crit}}{\pi}\;\frac{2^{\alpha -\frac{5}{2}} \Gamma (1-2 \alpha ) \left(\beta +\frac{t^2}{\beta }\right)^{\frac{1}{2}-\alpha } U\left(\frac{1}{2}-\alpha ,\frac{1}{2},\frac{2E_\mu \left(t^2+\beta ^2\right)}{\beta }\right)}{ \sqrt{\beta ^2+t^2}}.\label{eq:sff}
\end{equation}
This is shown in Figure \ref{fig:sff}, where we can observe that the spectral form factor grows linearly without reaching a plateau. Regardless, this is a computation at $g=0$ and $\hbar\rightarrow 0$, so the possibility that these and more non-perturbative corrections change the late time behavior or \eqref{eq:sff} cannot be ruled out. It is also worth noting that the low temperature and late time ($\beta\gg 1$ and $t\gg 0$) behavior of the spectral form factor does not depend on the values of $\alpha$.     
\begin{figure}[htbp]
\centering
\includegraphics[width=.9\textwidth]{spectraff.pdf}
\caption{Spectral form factor \eqref{eq:sff} for different values of $\alpha$, $t_{crit}=1$, $\beta=50$ and $E_\mu=0.1$.    \label{fig:sff}}
\end{figure}\\
\section{Discussion and outlook}\label{Section_Discussion}
Probably the most important result of this paper is to show how JT gravity and random geometry are deeply connected and how tools from both fields are useful for computational goals but also for our conceptual understanding. In this paper, we explored the random geometry aspects of JT gravity, offering insights into its critical regimes. 
Furthermore, we have remarked once again `the power of the string equation', already highlighted in multiple works \cite{Johnson:2019eik,Johnson:2021tnl,Turiaci:2020fjj}, noting its capacity not only to compute partition functions and density of states at the leading order but as a translation tool between random geometry and JT gravity. \\
By using tools coming from the study of Boltzmann planar maps, we uncovered what we call critical phases of JT gravity, which bear resemblance to phase transitions observed in planar maps, from surfaces with microscopic holes in the subcritical regime, to those with macroscopic ones in the non-generic critical regime. The occurrence of such transitions is the result of fine-tuning the weights of geodesic boundaries, which can be thought of as brane configurations. Furthermore, we made this analogy precise by determining the relation between hyperbolic surfaces with non-generic critical weights with $3/2<\alpha<5/2$ and the hyperbolic $O(n)$ loop model with $n=2\sin\left(\pi\left(\alpha-\frac{1}{2}\right)\right)$. We showed that the different values of $\alpha$ significantly impact the density of states of the system. Through our analysis, we have identified a family of models that smoothly interpolate between the density of states behaviors of $\rho_0(E)\sim\sqrt{E-E_0}$ and $\rho_0(E)\sim (E-E_0)^{3/2}$, thereby establishing connections to JT gravity coupled to dynamical EOW branes \cite{Gao:2021uro} and FZZT branes \cite{Okuyama:2021Page}. Moreover, we showed in which strict sense this transition from $\rho_0(E)\sim\sqrt{E-E_0}$ to $\rho_0(E)\sim (E-E_0)^{3/2}$ in the density of states is a proper phase transition in the physical sense and in which sense it is different compared to other transitions that have been identified previously \cite{Johnson:2020lns}. This is when the expected number of geodesic boundaries, or defects, diverge. Additionally, by analysing the effective dilaton potential for these models, one can distinguish between two subsets of models, $\alpha>2$ and $\alpha<2$ which suggests another phase transition at $\alpha=2$.\\
In light of our results, there are several interesting paths for future research to explore. 
Firstly, it remains essential to address the non-perturbative stability of the identified critical regimes and their compatibility with alternative definitions of EOW branes \cite{Johnson:2021tnl}.
This is related to an important question that arises, whether there exists a matrix model whose double-scaled density of states follows the form $(E-E_0)^{\alpha-1}$. Such a matrix model was constructed in \cite{Ambjorn:2016lkl}, but its non-perturbative study poses a big challenge due to the infinite order of its KdV differential equations \cite{Ambjorn:2017ibv}. We plan on addressing this  and investigating the properties of such a matrix model.
\\
Additionally, from the mathematical point of view, exploring the phase transition at $\alpha=2$ in the hyperbolic $O(n)$ loop model is an open problem. Our conjecture is that a transition from non-touching geodesic boundaries to overlapping ones should occur similarly to the planar map case.\\
Furthermore, we have presented a way to couple JT gravity to any kind of (admissible) defects using the string equation. This allows us to compare different results in the JT gravity literature in a single and precise context. For example, in \cite{Blommaert:2021fob}, the weights are chosen to achieve factorization in JT gravity. Such weight, $\mu(\ell_1,\ell_2)=-\frac{1}{\ell_1}\delta(\ell_1-\ell_2)$, is bilocal and negative valued which is not compatible with a probabilistic interpretation, but allows to `unglue' the double trumpet allowing for the factorization of the two boundary amplitude $\langle Z(\beta_1) Z(\beta_2)\rangle$. An even more interesting case for us is the weight choice used in \cite{Blommaert:2021gha} to obtain a discrete JT gravity spectrum. In this context, our non-generic critical weights are equivalent to finding a background matrix $H_0$ such that $\mathrm{Tr}(H_0)^k\sim k^{-\alpha}$ for $k\gg 1$. In this regime, when big holes appear in the hyperbolic surface, the density of states $\rho_0$ is conjectured to become discrete. Nevertheless, our results for $\rho_0$ do not agree with it, its leading order behavior changes from a square root in a non-trivial way but the spectrum stays continuum for all $\alpha$. The incompatibility of these two results poses an interesting avenue for further research. The last case we want to comment on is \cite{Okuyama:2021Page}, where the weights of anti-FZZT branes were tuned to reproduce a Page-like curve which exactly corresponds to a transition from $\rho_0(E)\sim\sqrt{E-E_0}$ to $\rho_0(E)\sim (E-E_0)^{3/2}$. Since our models interpolate between these two regimes, we think that the exploration of the non-generic critical phase raises intriguing questions regarding the implications for the Page curve and black hole evaporation.\\
Finally, we want to make a comment on the significance of non-generic critical models from the holographic point of view. One could forget about the geometry of the bulk and only look at the density of states of the quantum system at the holographic boundary. Then, in the subcritical regime $\rho_0(E)\sim\sqrt{E-E_0}$, which is the leading order behavior of the Schwarzian. However, in the non-generic critical regime, the density of states behaves as $\rho_0(E)\sim (E-E_0)^{\alpha-1}$, so a question that arises is if the transition from a subcritical to a non-generic critical regime can be put in the framework of deformations of the Schwarzian or if there are quantum systems that naturally present this the non-generic critical behavior, but we leave this for future work. 


\acknowledgments
I am grateful to Timothy Budd, Clifford Johnson, Bart Zonneveld and Pjotr Koster for valuable discussions and comments on a draft. In particular, Timothy Budd contributed to the analysis of the hyperbolic $O(n)$ loop model in Section \ref{section:Random_geometry}. This work is part of the START-UP 2018 programme with project number 740.018.017, which is financed by the Dutch Research Council (NWO).


% Bibliography

%% [A] Recommended: using JHEP.bst file
\bibliographystyle{JHEP}
\bibliography{biblio.bib}

\end{document}
